\documentclass[a4paper,10pt]{article}

\usepackage[utf8]{inputenc}
\usepackage[T1]{fontenc}
\usepackage[brazil]{babel}

\usepackage{geometry}
\usepackage{parskip}
\usepackage{microtype}
\usepackage{enumitem}
\usepackage{titlesec}
\usepackage{array}
\usepackage{tabularx}

\usepackage{hyperref}
\usepackage{bookmark}
\usepackage{xurl}

\geometry{top=1.5cm, bottom=1.5cm, left=1.5cm, right=1.5cm}
\setcounter{secnumdepth}{0}
\setlist[itemize]{
    leftmargin=0.75em,
    itemsep=0.2em,
    topsep=0.25em,
    parsep=0em,
    partopsep=0em
}
\pagestyle{empty}

\hypersetup{
    pdftitle={Currículo},
    colorlinks=true,
    linkcolor=black,
    urlcolor=black,
    citecolor=black,
    bookmarksdepth=2
}

\titleformat{\section}
{\Large\bfseries}
{}
{0em}
{}
[\titlerule\vspace{0.5ex}]
\titleformat{\subsection}
{\normalsize\bfseries}
{}
{0em}
{}
\titlespacing*{\subsection}{0pt}{0.35em}{0.15em}

\newcounter{cventry}
\newcommand{\cventry}[4]{%
    \refstepcounter{cventry}%
    \phantomsection%
    \pdfbookmark[2]{#1}{cventry-\thecventry}%
    \noindent\begin{tabularx}{\textwidth}{@{}>{\raggedright\arraybackslash}X >{\raggedleft\arraybackslash}X@{}}
    \textbf{#1} & #2 \\
    \textit{#3} & \textit{#4} \\
    \end{tabularx}
}

\begin{document}

\begin{center}
    {\LARGE \textbf{Lucas Ferreira Martins}} \\
    {\large Desenvolvedor de Software Júnior} \\
    São Paulo, SP
    {\textbullet}
    (11) 99999-1234
    {\textbullet}
    \href{mailto:lucas.martins@email.com}{lucas.martins@email.com} \\
    \url{linkedin.com/in/lucasferreiramartins}
    {\textbullet}
    \url{github.com/lucasferreira-dev}
\end{center}

\section*{Resumo}
Profissional de tecnologia com experiência no desenvolvimento de aplicações web e conhecimentos em JavaScript, React, Node.js e bancos de dados.

\section{Experiência}
\cventry{TechSolutions}{}{Desenvolvedor Web Júnior}{Março de 2024 – Atual}
\begin{itemize}
    \item Desenvolvimento e manutenção de páginas e funcionalidades de sistemas web.
    \item Criação de componentes utilizando React e JavaScript.
    \item Implementação de funcionalidades no servidor usando Node.js.
    \item Realização de consultas em bancos de dados.
    \item Correção de bugs e melhorias em aplicações existentes.
    \item Participação em reuniões de equipe e organização de tarefas.
\end{itemize}

\cventry{DigitalWeb}{}{Estagiário de Desenvolvimento}{Janeiro de 2023 – Fevereiro de 2024}
\begin{itemize}
    \item Auxílio no desenvolvimento de sites e sistemas internos.
    \item Utilização de HTML, CSS e JavaScript.
    \item Suporte na identificação e correção de problemas.
    \item Participação em testes de funcionalidades.
    \item Utilização do Git para controle de versões.
\end{itemize}

\section{Projetos}
\cventry{Sistema de Gerenciamento de Tarefas}{}{}{}
\begin{itemize}
    \item Desenvolvimento de aplicação para cadastro e gerenciamento de tarefas com React no frontend, Node.js no backend e PostgreSQL para armazenamento dos dados.
\end{itemize}

\cventry{Sistema de Cadastro de Clientes}{}{}{}
\begin{itemize}
    \item Criação de sistema web com funcionalidades de cadastro, edição, listagem e exclusão de clientes, utilizando JavaScript, Express e MySQL.
\end{itemize}

\section{Habilidades}
\begin{itemize}
    \item JavaScript
    \item TypeScript
    \item React
    \item Node.js
    \item HTML e CSS
    \item PostgreSQL
    \item MySQL
    \item Git e GitHub
    \item APIs REST
    \item Noções de Docker
    \item Metodologias ágeis
    \item Português: Nativo
    \item Inglês: Intermediário
\end{itemize}

\section{Educação}
\cventry{Universidade Paulista (UNIP)}{}{Bacharelado em Sistemas de Informação}{2022 – 2026}

\end{document}
